\documentclass[12pt]{article}
\usepackage[a4paper,margin=1in]{geometry}\usepackage[T1]{fontenc}
\usepackage{lmodern,microtype}\usepackage{amsmath,amssymb,amsthm,mathtools}
\usepackage{booktabs,array,enumitem,xcolor}\usepackage[numbers,sort&compress]{natbib}
\usepackage[colorlinks=true,linkcolor=blue!55!black,citecolor=blue!55!black,urlcolor=blue!55!black]{hyperref}
\linespread{1.07}
\newtheorem{theorem}{Theorem}[section]\newtheorem{proposition}[theorem]{Proposition}
\theoremstyle{remark}\newtheorem{remark}[theorem]{Remark}
\newcommand{\norm}[1]{\left\lVert#1\right\rVert}

\title{Dual-Channel Coherence of the Quartic Spectral Divisor\\
\large Newton Traces and a Nonstationary Small-Noise Flow}
\author{Liang Wang}\date{July 2026}
\begin{document}\maketitle
\begin{abstract}
State and residue transport are unstable under the naive adjacent-level map.
This paper tests the basis-invariant object that remains: the monic
characteristic polynomial of the source-observable edge quartet.

At each of $\sigma=0.04,0.02,0.01$, the left and right physical channels
produce quartic coefficient vectors agreeing to relative $\ell^2$ error at
most $0.008112$.  Their constant terms agree to relative errors
$0.02716$, $0.01301$, and $0.007749$, respectively.  This channel coherence
is substantially stronger than the order-one projector and residue
transport defects.

Across scales the divisor is not stationary: adjacent coefficient errors
range from $0.24023$ to $0.31674$.  We prove that its coefficients and all
unweighted power traces are equivalent through Newton identities and verify
the reconstruction through power ten on 120 random complex quartets with
zero failures.  The result identifies a promising scalar flow, not a limit:
three levels do not establish coefficient compactness, a renormalization
law, a Fredholm determinant, or Gate A.
\end{abstract}

\section{Why pass from states to a divisor}

RH-202 rejects the direct Haar transport of the outer eigenspaces
\citep{WangRH202}, and RH-206 rejects one common scalar for their physical
residues \citep{WangRH206}.  Both failures concern coordinate- or
source-dependent objects.  The multiset of quartet eigenvalues is invariant
under similarity and independent of eigenvector normalization.

For an ordered or unordered quartet $\Lambda=\{\lambda_1,\ldots,\lambda_4\}$
define
\begin{equation}\label{eq:divisor}
 D_\Lambda(z)=\prod_{j=1}^4(z-\lambda_j)
 =z^4+c_1z^3+c_2z^2+c_3z+c_4.
\end{equation}
For two conjugate pairs all $c_j$ are real in exact arithmetic.  We call the
coefficient vector the finite spectral divisor ledger.

\section{Similarity and permutation invariance}

\begin{proposition}[Finite divisor invariance]\label{prop:invariance}
The coefficients in \eqref{eq:divisor} are invariant under permutation of
the quartet and under similarity of any matrix whose restricted spectrum is
$\Lambda$.  The constant term is
\begin{equation}\label{eq:det}
 c_4=\prod_{j=1}^4\lambda_j=\det K_\Lambda.
\end{equation}
\end{proposition}

This allows the left and right physical channels to be compared without
constructing a common eigenvector map.

\section{Newton trace equivalence}

Define the unweighted power sums
\begin{equation}\label{eq:power}
 p_q=\sum_{j=1}^4\lambda_j^q.
\end{equation}

\begin{theorem}[Newton recursion]\label{thm:newton}
For $1\le q\le4$,
\begin{equation}\label{eq:newton-low}
 p_q+c_1p_{q-1}+\cdots+c_{q-1}p_1+qc_q=0.
\end{equation}
For $q>4$,
\begin{equation}\label{eq:newton-high}
 p_q+c_1p_{q-1}+c_2p_{q-2}+c_3p_{q-3}+c_4p_{q-4}=0.
\end{equation}
Thus the monic quartic determines every finite power trace, and
$p_1,\ldots,p_4$ determine the quartic.
\end{theorem}

\begin{proof}
These are the Newton--Girard identities applied to the elementary symmetric
functions of the roots.  The first four equations solve successively for
$c_1,\ldots,c_4$; the characteristic equation gives the later recurrence.
\end{proof}

This ledger is unweighted.  It must not be confused with physical moments
$\sum_jr_j\lambda_j^q$, whose residues obey the RH-206 cocycle.

\section{Three physical quartics}

The real coefficient vectors are:
\begin{center}
\small
\begin{tabular}{llrrrr}
\toprule
$\sigma$ & side & $c_1$ & $c_2$ & $c_3$ & $c_4$\\
\midrule
$0.04$ & left  & $-0.06975$ & $0.56039$ & $-0.07015$ & $0.05860$\\
$0.04$ & right & $-0.07789$ & $0.55866$ & $-0.07409$ & $0.05701$\\
$0.02$ & left  & $ 0.06254$ & $0.55984$ & $-0.02530$ & $0.30762$\\
$0.02$ & right & $ 0.06128$ & $0.55822$ & $-0.02615$ & $0.30362$\\
$0.01$ & left  & $ 0.04238$ & $0.22166$ & $-0.06578$ & $0.33826$\\
$0.01$ & right & $ 0.04295$ & $0.22785$ & $-0.06733$ & $0.33564$\\
\bottomrule
\end{tabular}
\end{center}
Imaginary coefficient remnants are at roundoff scale because every packet
is conjugation closed.

\section{Left/right channel coherence}

For coefficient vectors $c^L,c^R$, record
\begin{equation}\label{eq:error}
 \varepsilon_{LR}=\frac{\norm{c^L-c^R}_2}{\norm{c^L}_2}.
\end{equation}
The three errors are $0.008112$, $0.003851$, and $0.006398$.  All lie below
one percent.

The constant term discrepancies decrease across the three anchors:
\begin{equation}\label{eq:constant-flow}
 0.02716,\qquad0.01301,\qquad0.007749.
\end{equation}
This is a descriptive finite trend.  It is consistent with the two channels
approximating one scalar spectral object, but does not prove that their
difference tends to zero.

\section{Scale flow is not yet converged}

Adjacent coefficient-vector errors are:
\begin{center}
\begin{tabular}{lrr}
\toprule
step & left & right\\
\midrule
$0.04\to0.02$ & $0.24023$ & $0.24202$\\
$0.02\to0.01$ & $0.31674$ & $0.30955$\\
\bottomrule
\end{tabular}
\end{center}
Unlike the left/right discrepancy, the scale movement is not small and does
not decrease on the finer transition.  A raw coefficient limit cannot be
inferred from these points.

The coefficient changes also occur in different coordinates: $c_4$ rises
strongly from $\sigma=0.04$ to $0.02$, while $c_2$ changes strongly from
$0.02$ to $0.01$.  Any renormalization should therefore be tested on the
full vector rather than fitted to one coefficient.

\section{Two intrinsic normalization candidates}

The divisor itself supplies scale and center statistics.  First define the
determinant radius
\begin{equation}\label{eq:rho}
 \rho=|c_4|^{1/4}.
\end{equation}
Under $\lambda_j'=\lambda_j/\rho$, the normalized coefficients are
\begin{equation}\label{eq:radial-coeff}
 c_j'=c_j/\rho^j,
\end{equation}
so no levelwise regression parameter is introduced.

A second candidate removes translation.  Since
$\sum_j\lambda_j=-c_1$, the intrinsic center is
\begin{equation}\label{eq:center}
 \mu=-c_1/4.
\end{equation}
Translate $\lambda_j\mapsto\lambda_j-\mu$ and divide by the centered root
mean square radius.  The resulting cubic coefficient vanishes exactly, and
the quadratic energy is normalized.  These two formulas are predeclared
candidates for RH-212; neither is claimed to improve the present data until
tested on denser anchors.

\section{Machine audit of Newton equivalence}

For 120 random complex quartets we form the monic polynomial, reconstruct
$p_1,\ldots,p_{10}$ with \eqref{eq:newton-low}--\eqref{eq:newton-high}, and
compare with direct sums of powers.  The maximum floating error is below the
declared $10^{-8}$ gate and all cases pass.  This verifies that the archived
coefficient and trace ledgers use one exact algebraic convention.

\section{A scalar route and its requirements}

The finite evidence suggests studying a map
\begin{equation}\label{eq:flow}
 c(\sigma)\longmapsto c(\sigma/2)
\end{equation}
before insisting on convergence of raw projectors.  A successful Gate-A
argument would need at least:
\begin{enumerate}[leftmargin=2em]
\item denser small-noise anchors and a predeclared branch rule;
\item compactness or a contraction after intrinsic coefficient
renormalization;
\item control of omitted modes as the packet grows;
\item convergence of finite determinants on a common complex domain;
\item only then, a Fredholm or dynamical determinant interpretation.
\end{enumerate}
The present quartic is a local test block and can never by itself produce an
infinite spectral counting law.

\section{From one quartic to a determinant family}

Suppose later contours select packets $\Lambda_{\sigma,k}$ of growing size.
A meaningful determinant limit requires more than coefficientwise
convergence at fixed $k$: the finite products must form a normal family on a
common domain, and tails must be controlled uniformly in both $\sigma$ and
$k$.  Canonical products or regularized determinants may be required when
ordinary products diverge \citep{Simon2005}.  RH-207 supplies only the
$k=4$ laboratory in which candidate normalizations can be falsified cheaply.

\section{What is not encoded}

The divisor contains eigenvalues with multiplicity, but not the physical
source residues, projector conditioning, or eigenvectors.  It has no
established arithmetic weights.  In particular, the Newton traces here are
not von Mangoldt traces, and the roots are not identified with zeta zeros.

\section{Claim boundary and next step}

Similarity invariance and Newton equivalence are exact.  Channel coherence
and scale movement are finite floating observations at three levels.  No
coefficient limit, Fredholm determinant, $T\log T$ count, Gate A, or later
macro gate is claimed.  RH-208 next separates endpoint spectral isolation
from the much harder certification of interlevel transport.

\bibliographystyle{plainnat}\bibliography{references}\end{document}
